\documentclass[12pt,a4paper]{article}
\usepackage[UTF8]{ctex}
\usepackage{amsmath,amssymb,amsthm}
\usepackage{geometry}
\usepackage{bm}

\geometry{margin=2.5cm}

\title{动能公式中$V_b^T$的展开：$\frac{1}{2}V_b^T G_b V_b$}
\author{第8章补充说明}
\date{}

\begin{document}

\maketitle

\section{问题提出}

\textbf{问题}：动能公式
\begin{equation}
K = \frac{1}{2}\omega_b^T I_b \omega_b + \frac{1}{2}mv_b^T v_b = \frac{1}{2}V_b^T G_b V_b
\label{eq:kinetic_energy}
\end{equation}

中的$V_b^T$是什么？需要详细展开。

\section{符号定义}

\subsection{Body Twist $V_b$}

\textbf{定义}：Body twist $V_b$是一个6维向量，包含角速度和线速度：
\begin{equation}
V_b = \begin{bmatrix} \omega_b \\ v_b \end{bmatrix} \in \mathbb{R}^6
\label{eq:Vb_definition}
\end{equation}

其中：
\begin{itemize}
\item $\omega_b \in \mathbb{R}^3$：角速度（在体坐标系$\{b\}$中表示）
\item $v_b \in \mathbb{R}^3$：线速度（在体坐标系$\{b\}$中表示）
\end{itemize}

\subsection{$V_b^T$的展开}

\textbf{转置}：
\begin{equation}
V_b^T = \begin{bmatrix} \omega_b \\ v_b \end{bmatrix}^T = \begin{bmatrix} \omega_b^T & v_b^T \end{bmatrix} \in \mathbb{R}^{1 \times 6}
\label{eq:Vb_transpose}
\end{equation}

\textbf{详细展开}：

设$\omega_b = \begin{bmatrix} \omega_x \\ \omega_y \\ \omega_z \end{bmatrix}$，$v_b = \begin{bmatrix} v_x \\ v_y \\ v_z \end{bmatrix}$，则：
\begin{equation}
V_b^T = \begin{bmatrix} \omega_x & \omega_y & \omega_z & v_x & v_y & v_z \end{bmatrix}
\end{equation}

\subsection{空间惯性矩阵$G_b$}

\textbf{定义}：
\begin{equation}
G_b = \begin{bmatrix} I_b & 0 \\ 0 & mI \end{bmatrix} \in \mathbb{R}^{6 \times 6}
\label{eq:Gb_definition}
\end{equation}

其中：
\begin{itemize}
\item $I_b \in \mathbb{R}^{3 \times 3}$：转动惯量矩阵
\item $m \in \mathbb{R}$：刚体的质量
\item $I \in \mathbb{R}^{3 \times 3}$：$3 \times 3$单位矩阵
\item $0 \in \mathbb{R}^{3 \times 3}$：$3 \times 3$零矩阵
\end{itemize}

\textbf{详细展开}：
\begin{equation}
G_b = \begin{bmatrix}
I_{xx} & I_{xy} & I_{xz} & 0 & 0 & 0 \\
I_{xy} & I_{yy} & I_{yz} & 0 & 0 & 0 \\
I_{xz} & I_{yz} & I_{zz} & 0 & 0 & 0 \\
0 & 0 & 0 & m & 0 & 0 \\
0 & 0 & 0 & 0 & m & 0 \\
0 & 0 & 0 & 0 & 0 & m
\end{bmatrix}
\end{equation}

\section{动能公式的展开}

\subsection{方法1：直接展开$V_b^T G_b V_b$}

\textbf{步骤1}：计算$G_b V_b$

\begin{align}
G_b V_b &= \begin{bmatrix} I_b & 0 \\ 0 & mI \end{bmatrix} \begin{bmatrix} \omega_b \\ v_b \end{bmatrix} \\
&= \begin{bmatrix} I_b \omega_b \\ mI v_b \end{bmatrix} \\
&= \begin{bmatrix} I_b \omega_b \\ mv_b \end{bmatrix}
\label{eq:Gb_Vb}
\end{align}

\textbf{步骤2}：计算$V_b^T (G_b V_b)$

\begin{align}
V_b^T G_b V_b &= \begin{bmatrix} \omega_b^T & v_b^T \end{bmatrix} \begin{bmatrix} I_b \omega_b \\ mv_b \end{bmatrix} \\
&= \omega_b^T I_b \omega_b + v_b^T (mv_b) \\
&= \omega_b^T I_b \omega_b + m v_b^T v_b
\label{eq:VbT_Gb_Vb}
\end{align}

\textbf{步骤3}：动能公式

\begin{equation}
K = \frac{1}{2}V_b^T G_b V_b = \frac{1}{2}\omega_b^T I_b \omega_b + \frac{1}{2}mv_b^T v_b
\end{equation}

\textbf{证毕}。

\subsection{方法2：分块矩阵乘法}

\textbf{使用分块矩阵}：

\begin{align}
V_b^T G_b V_b &= \begin{bmatrix} \omega_b^T & v_b^T \end{bmatrix} \begin{bmatrix} I_b & 0 \\ 0 & mI \end{bmatrix} \begin{bmatrix} \omega_b \\ v_b \end{bmatrix} \\
&= \begin{bmatrix} \omega_b^T & v_b^T \end{bmatrix} \begin{bmatrix} I_b \omega_b \\ mI v_b \end{bmatrix} \\
&= \omega_b^T I_b \omega_b + v_b^T (mI) v_b \\
&= \omega_b^T I_b \omega_b + m v_b^T v_b
\end{align}

\subsection{方法3：逐元素展开}

\textbf{详细展开每个分量}：

\textbf{第一部分}：$\omega_b^T I_b \omega_b$

设$\omega_b = \begin{bmatrix} \omega_x \\ \omega_y \\ \omega_z \end{bmatrix}$，$I_b = \begin{bmatrix}
I_{xx} & I_{xy} & I_{xz} \\
I_{xy} & I_{yy} & I_{yz} \\
I_{xz} & I_{yz} & I_{zz}
\end{bmatrix}$，则：

\begin{align}
\omega_b^T I_b \omega_b &= \begin{bmatrix} \omega_x & \omega_y & \omega_z \end{bmatrix} \begin{bmatrix}
I_{xx} & I_{xy} & I_{xz} \\
I_{xy} & I_{yy} & I_{yz} \\
I_{xz} & I_{yz} & I_{zz}
\end{bmatrix} \begin{bmatrix} \omega_x \\ \omega_y \\ \omega_z \end{bmatrix} \\
&= \begin{bmatrix} \omega_x & \omega_y & \omega_z \end{bmatrix} \begin{bmatrix}
I_{xx}\omega_x + I_{xy}\omega_y + I_{xz}\omega_z \\
I_{xy}\omega_x + I_{yy}\omega_y + I_{yz}\omega_z \\
I_{xz}\omega_x + I_{yz}\omega_y + I_{zz}\omega_z
\end{bmatrix} \\
&= I_{xx}\omega_x^2 + I_{yy}\omega_y^2 + I_{zz}\omega_z^2 + 2I_{xy}\omega_x\omega_y + 2I_{xz}\omega_x\omega_z + 2I_{yz}\omega_y\omega_z
\end{align}

\textbf{第二部分}：$m v_b^T v_b$

设$v_b = \begin{bmatrix} v_x \\ v_y \\ v_z \end{bmatrix}$，则：

\begin{align}
m v_b^T v_b &= m \begin{bmatrix} v_x & v_y & v_z \end{bmatrix} \begin{bmatrix} v_x \\ v_y \\ v_z \end{bmatrix} \\
&= m(v_x^2 + v_y^2 + v_z^2) \\
&= m \|v_b\|^2
\end{align}

\textbf{总动能}：

\begin{equation}
K = \frac{1}{2}\left[I_{xx}\omega_x^2 + I_{yy}\omega_y^2 + I_{zz}\omega_z^2 + 2I_{xy}\omega_x\omega_y + 2I_{xz}\omega_x\omega_z + 2I_{yz}\omega_y\omega_z + m(v_x^2 + v_y^2 + v_z^2)\right]
\end{equation}

\section{物理意义}

\subsection{旋转动能}

\textbf{第一项}：$\frac{1}{2}\omega_b^T I_b \omega_b$

这是刚体的\textbf{旋转动能}，表示由于旋转运动而具有的能量。

\textbf{展开形式}：
\begin{equation}
K_{\text{rot}} = \frac{1}{2}(I_{xx}\omega_x^2 + I_{yy}\omega_y^2 + I_{zz}\omega_z^2 + 2I_{xy}\omega_x\omega_y + 2I_{xz}\omega_x\omega_z + 2I_{yz}\omega_y\omega_z)
\end{equation}

\textbf{物理意义}：
\begin{itemize}
\item 对角项（$I_{xx}\omega_x^2$等）：绕各轴的旋转动能
\item 交叉项（$2I_{xy}\omega_x\omega_y$等）：由于惯性积产生的耦合项
\end{itemize}

\subsection{平移动能}

\textbf{第二项}：$\frac{1}{2}mv_b^T v_b$

这是刚体的\textbf{平移动能}，表示由于平移运动而具有的能量。

\textbf{展开形式}：
\begin{equation}
K_{\text{trans}} = \frac{1}{2}m(v_x^2 + v_y^2 + v_z^2) = \frac{1}{2}m\|v_b\|^2
\end{equation}

\textbf{物理意义}：这是质心的动能，等于$\frac{1}{2}mv^2$，其中$v = \|v_b\|$是质心速度的大小。

\subsection{总动能}

\textbf{总动能}：
\begin{equation}
K = K_{\text{rot}} + K_{\text{trans}} = \frac{1}{2}\omega_b^T I_b \omega_b + \frac{1}{2}mv_b^T v_b
\end{equation}

\textbf{物理意义}：刚体的总动能等于旋转动能和平移动能之和。

\section{矩阵乘法的详细步骤}

\subsection{步骤1：$V_b^T$的定义}

\begin{equation}
V_b^T = \begin{bmatrix} \omega_b^T & v_b^T \end{bmatrix} = \begin{bmatrix} \omega_x & \omega_y & \omega_z & v_x & v_y & v_z \end{bmatrix}
\end{equation}

这是一个$1 \times 6$的行向量。

\subsection{步骤2：$G_b$的定义}

\begin{equation}
G_b = \begin{bmatrix}
I_b & 0 \\
0 & mI
\end{bmatrix} = \begin{bmatrix}
I_{xx} & I_{xy} & I_{xz} & 0 & 0 & 0 \\
I_{xy} & I_{yy} & I_{yz} & 0 & 0 & 0 \\
I_{xz} & I_{yz} & I_{zz} & 0 & 0 & 0 \\
0 & 0 & 0 & m & 0 & 0 \\
0 & 0 & 0 & 0 & m & 0 \\
0 & 0 & 0 & 0 & 0 & m
\end{bmatrix}
\end{equation}

这是一个$6 \times 6$的矩阵。

\subsection{步骤3：$G_b V_b$的计算}

\begin{align}
G_b V_b &= \begin{bmatrix}
I_b & 0 \\
0 & mI
\end{bmatrix} \begin{bmatrix} \omega_b \\ v_b \end{bmatrix} \\
&= \begin{bmatrix}
I_b \omega_b \\
mI v_b
\end{bmatrix} \\
&= \begin{bmatrix}
I_{xx}\omega_x + I_{xy}\omega_y + I_{xz}\omega_z \\
I_{xy}\omega_x + I_{yy}\omega_y + I_{yz}\omega_z \\
I_{xz}\omega_x + I_{yz}\omega_y + I_{zz}\omega_z \\
mv_x \\
mv_y \\
mv_z
\end{bmatrix}
\end{align}

这是一个$6 \times 1$的列向量。

\subsection{步骤4：$V_b^T (G_b V_b)$的计算}

\begin{align}
V_b^T G_b V_b &= \begin{bmatrix} \omega_x & \omega_y & \omega_z & v_x & v_y & v_z \end{bmatrix} \begin{bmatrix}
I_{xx}\omega_x + I_{xy}\omega_y + I_{xz}\omega_z \\
I_{xy}\omega_x + I_{yy}\omega_y + I_{yz}\omega_z \\
I_{xz}\omega_x + I_{yz}\omega_y + I_{zz}\omega_z \\
mv_x \\
mv_y \\
mv_z
\end{bmatrix} \\
&= \omega_x(I_{xx}\omega_x + I_{xy}\omega_y + I_{xz}\omega_z) \\
&\quad + \omega_y(I_{xy}\omega_x + I_{yy}\omega_y + I_{yz}\omega_z) \\
&\quad + \omega_z(I_{xz}\omega_x + I_{yz}\omega_y + I_{zz}\omega_z) \\
&\quad + v_x(mv_x) + v_y(mv_y) + v_z(mv_z) \\
&= I_{xx}\omega_x^2 + I_{xy}\omega_x\omega_y + I_{xz}\omega_x\omega_z \\
&\quad + I_{xy}\omega_x\omega_y + I_{yy}\omega_y^2 + I_{yz}\omega_y\omega_z \\
&\quad + I_{xz}\omega_x\omega_z + I_{yz}\omega_y\omega_z + I_{zz}\omega_z^2 \\
&\quad + m(v_x^2 + v_y^2 + v_z^2) \\
&= I_{xx}\omega_x^2 + I_{yy}\omega_y^2 + I_{zz}\omega_z^2 \\
&\quad + 2I_{xy}\omega_x\omega_y + 2I_{xz}\omega_x\omega_z + 2I_{yz}\omega_y\omega_z \\
&\quad + m(v_x^2 + v_y^2 + v_z^2)
\end{align}

\textbf{因此}：
\begin{equation}
K = \frac{1}{2}V_b^T G_b V_b = \frac{1}{2}\omega_b^T I_b \omega_b + \frac{1}{2}mv_b^T v_b
\end{equation}

\section{特殊情况}

\subsection{主惯性轴情况}

如果坐标轴与主惯性轴对齐，则$I_b$是对角矩阵：
\begin{equation}
I_b = \begin{bmatrix}
I_{xx} & 0 & 0 \\
0 & I_{yy} & 0 \\
0 & 0 & I_{zz}
\end{bmatrix}
\end{equation}

\textbf{旋转动能简化为}：
\begin{equation}
K_{\text{rot}} = \frac{1}{2}(I_{xx}\omega_x^2 + I_{yy}\omega_y^2 + I_{zz}\omega_z^2)
\end{equation}

没有交叉项。

\subsection{纯旋转情况}

如果$v_b = 0$（只有旋转，没有平移），则：
\begin{equation}
K = \frac{1}{2}\omega_b^T I_b \omega_b
\end{equation}

\subsection{纯平移情况}

如果$\omega_b = 0$（只有平移，没有旋转），则：
\begin{equation}
K = \frac{1}{2}mv_b^T v_b = \frac{1}{2}m\|v_b\|^2
\end{equation}

\section{总结}

\subsection{主要结论}

\begin{enumerate}
\item \textbf{$V_b^T$的定义}：
\begin{equation}
V_b^T = \begin{bmatrix} \omega_b^T & v_b^T \end{bmatrix} = \begin{bmatrix} \omega_x & \omega_y & \omega_z & v_x & v_y & v_z \end{bmatrix}
\end{equation}

\item \textbf{动能公式的展开}：
\begin{equation}
K = \frac{1}{2}V_b^T G_b V_b = \frac{1}{2}\omega_b^T I_b \omega_b + \frac{1}{2}mv_b^T v_b
\end{equation}

\item \textbf{详细形式}：
\begin{align}
K &= \frac{1}{2}[I_{xx}\omega_x^2 + I_{yy}\omega_y^2 + I_{zz}\omega_z^2 \\
&\quad + 2I_{xy}\omega_x\omega_y + 2I_{xz}\omega_x\omega_z + 2I_{yz}\omega_y\omega_z \\
&\quad + m(v_x^2 + v_y^2 + v_z^2)]
\end{align}
\end{enumerate}

\subsection{关键要点}

\begin{itemize}
\item $V_b^T$是$1 \times 6$的行向量
\item $G_b$是$6 \times 6$的对称正定矩阵
\item $V_b^T G_b V_b$是标量（动能）
\item 动能由旋转动能和平移动能两部分组成
\end{itemize}

\subsection{应用场景}

\begin{itemize}
\item 刚体动力学分析
\item 机器人能量计算
\item 航天器姿态动力学
\item 机械系统能量分析
\end{itemize}

\end{document}

